\chapter{Evaluation}\label{chap:evaluation}

This chapter evaluates the lattice-based RISC-V ZkVM prototype developed in
this work. The current implementation is incomplete and should not be understood
as a production-ready ZkVM. The goal of this chapter is therefore not to claim
state-of-the-art performance, but to compare the prototype with existing ZkVMs,
to explain the observed differences, and to answer the research questions of
this thesis.

The performance measurements were conducted on server-grade hardware with the
following configuration:

\begin{itemize}
	\item \textbf{CPU:} Intel Xeon E312xx
	\item \textbf{OS:} Ubuntu 24.04
	\item \textbf{Memory:} 64 GB RAM
\end{itemize}

\section{Fibonacci benchmark}

A simple benchmark was conducted using a program that computes the 100th
Fibonacci number. This benchmark was selected because it is small, easy to
reproduce, and does not require a large number of RISC-V instructions.

The Fibonacci program execution in the prototype resulted in:

\begin{itemize}
	\item \textbf{Execution traces:} 16 instructions
	\item \textbf{Memory operations:} 2 operations
	\item \textbf{Total folding time:} 688.24 seconds
\end{itemize}

This result already shows that the current implementation is considerably slower
than modern ZkVM systems. However, this number must be interpreted carefully,
because the prototype still contains many unoptimized components and runs only
on CPU.

\section{Comparison with other ZkVMs}

Additional benchmark artifacts were collected on the same machine for Pico,
SP1, ZKsync Airbender, and Nexus zkVM 3.0. Table~\ref{tab:zkvm-benchmarks-v2}
compares the wall-clock time of the proving and the peak RAM usage.

ZisK is not included in this comparison. The benchmark server did not support
the SIMD instruction set required by ZisK, and therefore a comparable proof run
could not be executed on the same hardware.

\begin{table}[h]
	\centering
	\begin{tabular}{|l|c|c|}
		\hline
		\textbf{ZkVM}    & \textbf{Proof time} & \textbf{Peak RAM} \\
		\hline
		This work        & 11:28.92            & 12.34 GB          \\
		Pico             & 2:04.50             & 6.27 GB           \\
		SP1              & 1:50.17             & 9.96 GB           \\
		ZKsync Airbender & 4:06.73             & 24.34 GB          \\
		Nexus zkVM 3.0   & 0:22.81             & 0.12 GB           \\
		\hline
	\end{tabular}
	\caption{Comparison of proof generation time and peak RAM usage for the Fibonacci benchmark.}
	\label{tab:zkvm-benchmarks-v2}
\end{table}

The comparison shows that the proposed system is slower than the compared
systems. Pico and SP1 complete the proof much faster, Nexus zkVM 3.0 is also
substantially faster, and Airbender is also significantly faster in time,
although it uses substantially more memory.
As recent systematization work emphasizes, such differences are influenced not
only by the proof system itself, but also by ISA choice, trace layout, memory
checking strategy, and recursion architecture \cite{SoKZkVMs}.

At the same time, one possible advantage of the lattice-based construction may
be proof size. The current prototype does not yet produce the final Greyhound
wrapping proof, so no measured proof size can be reported at this
stage. Therefore, the following comparison is only theoretical.

For the compared ZkVMs, the reported final proof sizes are approximately 1 MB
for Airbender, approximately 1 MB for Pico, and approximately 260 KB for ZisK,
as summarized in the earlier ZkEVM comparison in this thesis. In contrast, the
Greyhound paper reports polynomial evaluation proofs of about 53 KB for large
polynomials of degree up to $N = 2^{30}$. This is not directly the same as a
final wrapping proof for the proposed ZkVM, so these values are not fully
comparable. However, they indicate that a future Greyhound-based wrapper could
potentially produce proofs in the order of tens of kilobytes. If this can be
achieved in a complete implementation, proof size could become an important
practical advantage of the lattice-based approach.

To complement the small Fibonacci benchmark, a longer experiment was also run
with the \texttt{evm} guest based on \texttt{revm}. This run continued for
approximately 9 hours and 29 minutes on the same machine. During this time, the
system kept processing the execution and verifying the IVC chain, while peak RAM
usage remained around 13.7 GB. This is still far from production-level
performance, but it is an encouraging sign that the implementation can continue
working on a substantially more complex workload without an uncontrolled memory
increase or failing to produce proof.

\section{Why is the prototype slower?}

The benchmark results suggest that the main limitation of the current prototype
is not correctness, but performance. Several reasons can explain this outcome.

\paragraph{Folding is not as cheap in practice as it appears in theory.}

The literature often describes folding as a sequence of relatively cheap random
linear combinations. From a theoretical perspective this description is correct.
However, the practical implementation cost is still high.

In the prototype, a significant amount of time is spent in multilinear
extension-related computations. Even if the algebraic step is conceptually only
a random linear combination, the real implementation still has to evaluate large
objects, move large vectors through memory, and repeatedly perform arithmetic on
high-dimensional representations. Therefore, even though LatticeFold is
efficient in theory, it is not automatically efficient in practice.

This observation is supported by profiling of the Fibonacci proving run. The
flame graph (which can be found in the project source code named\\\texttt{fibonacci\_flamegraph.svg}) shows that most of the
runtime is spent inside Rayon worker threads, where the dominant application
hotspots belong to \texttt{latticefold}'s multilinear extension and
folding/decomposition routines. The most latticefold related functions are
\texttt{compute\_mz\_mles}, \texttt{to\_mles\_err},
\texttt{evaluate\_mles}, \texttt{get\_etas},
\texttt{compute\_u\_s} and there is also visible time spent in
\texttt{DenseMultilinearExtension} arithmetic.

By contrast, the commitment computation is present in the profile but is not one
of the main bottlenecks, and witness arithmetization itself is not especially
prominent. The profile therefore suggests that the main practical cost of the
current prototype is not the front-end VM execution, but the MLE-heavy algebraic
work required by decomposition and folding. This is consistent with the measured
runtime, where the proving overhead is dominated by the algebraic representation
of the witness rather than by the small Fibonacci program itself.

\paragraph{CCS may be an inconvenient computation model for this use case.}

Another possible reason is the choice of CCS as the computation model.

CCS is flexible and is well suited for expressing linear relations. This makes
it a reasonable model for several simple VM instructions. However, the full
system does not contain only linear computation. The non-interactive folding
scheme itself is not purely linear, and the verifier logic introduces higher
degree structure.

This creates several problems in practice:

\begin{itemize}
	\item CCS does not provide custom gates in the same sense as many modern proving
	      systems. Because of that, the representation may require more matrices and
	      more bookkeeping than a more specialized system.
	\item Some parts of the VM are relatively simple in CCS, but components such as
	      Poseidon2 and the folding verifier are much less natural to express.
	\item The folding verifier is especially difficult because its constraints must be
	      self-aware. In other words, they must already know the shape of the proof
	      that they are verifying.
\end{itemize}

This last point may create a kind of circular dependency. To verify an object of
higher degree $D$, the verifier may require information whose own constraint
system increases the degree again. Then one has to enlarge the representation,
which may increase the maximal degree once more. In this sense, the model may
create a loop where expressing the verifier makes the verifier itself more
expensive to express.

This also suggests that lookup arguments may be a useful complement to the
current design. Lookup systems such as Lasso \cite{Lasso} can encode repeated
table relations more compactly than direct constraint expansion, which may
be relevant for operations on bytes, value checks, and other common VM subroutines.

For this reason, CCS may be acceptable for basic linear VM instructions, but it
may be a poor fit for cryptographic permutations such as Poseidon2 and for the
self-referential structure of folding verification.

\section{Potential speed improvements}

The current benchmark should not be understood as a final performance limit.
Several optimization directions remain open.

\subsection{SIMD and GPU acceleration}

The present implementation runs only on CPU. This immediately puts it at a
disadvantage compared to systems that use heavily optimized parallel proving
pipelines.

The arithmetic in the prototype has regular structure, which makes it suitable
for SIMD acceleration. Large vector operations, MLE evaluations, and some parts
of the commitment computations could benefit from parallel execution.
GPU acceleration is another natural direction, especially for large batches of
similar field or ring operations.

There is also some encouraging ecosystem progress in this direction. Recently,
a pull request was merged into the Arkworks algebra libraries whose primary
motivation was to create a generalized path toward vectorized SIMD
optimizations for finite fields in Arkworks\footnote{\url{https://github.com/arkworks-rs/algebra/pull/1044}.}.
This is relevant, because the current implementation depends on Arkworks-based components.
However, using this new path in practice would still require non-trivial work.
In particular, the \texttt{latticefold} library would need to rewrite important
parts of its internal representation to use the new SIMD-friendly field types.
Therefore, this optimization direction is promising, but it is not something
that can be enabled by a small configuration change.

\subsection{More mathematical optimizations}

Because the system is built over algebraic structures, it may admit more
specialized optimizations than a direct implementation currently uses.

Possible examples include:

\begin{itemize}
	\item more efficient layouts for multilinear extensions,
	\item better batching of linear algebra operations,
	\item reuse of intermediate evaluations across folding steps,
	\item specialization of constraints for the most common instruction classes.
\end{itemize}

These improvements are not only low-level engineering optimizations. They are
also mathematical optimizations, because the efficiency of the system depends
strongly on how the algebraic objects are represented and evaluated.

\section{Addressing the research questions}

\subsection*{Feasibility of lattice-based instantiation}

\textit{Is it feasible to instantiate a full RISC-V execution environment using
	CCS compatible with the LatticeFold protocol?}

\textbf{Answer:} Partially yes.

From the design perspective, the thesis demonstrates that such an architecture
can be specified. The mapping from RISC-V execution state to CCS-compatible
constraints was designed, and the connection to LatticeFold was established.

From the implementation perspective, the feasibility is only partially shown.
The prototype contains a functional RISC-V emulator, witness generation, and the
basic integration with folding. However, it still lacks complete instruction set
coverage and a complete in-circuit verifier for the full recursive workflow.

Therefore, the answer is positive on the architectural level, but still
incomplete on the implementation level.

\subsection*{Performance trade-offs}

\textit{How does the performance of a lattice-based folding scheme compare to
	traditional elliptic curve-based folding schemes when applied to general-purpose
	VM execution?}

\textbf{Answer:} In the current prototype, the lattice-based approach is slower.

The benchmark comparison shows that the implemented system needs more time than
the compared ZkVMs to generate a proof for the benchmark program. The main
practical reason appears to be that the algebraic work of folding, especially
MLE-related computation, is expensive in the current implementation. Another
reason is that CCS may not be the most convenient model for all parts of the
system, especially for cryptographic permutations and folding verification.

At the same time, the results do not imply that lattice-based folding is a bad
direction in principle. The implementation is still early, CPU-only, and not yet
optimized. In addition, compact proof size remains a positive property of the
approach.

Therefore, the present answer is that the current implementation has worse
practical proving time, but the long-term trade-off remains open.

\subsection*{Suitability for long-running computations}

\textit{Does the noise growth inherent in lattice-based cryptography hinder the
	ability to prove long execution traces, such as those required for Ethereum block
	validation?}

\textbf{Answer:} The current results are not conclusive, but they provide a
positive outlook.

From the design perspective, LatticeFold addresses witness norm growth by
decomposition during folding. This suggests that long computations should remain
manageable in principle.

From the implementation perspective, the Fibonacci benchmark is too small to
provide a strong answer, because 16 traces do not meaningfully stress the
long-run behavior of the system. However, the additional \texttt{revm}-based EVM
experiment gives a more encouraging signal. That run continued for more than 9
hours, the IVC chain kept being verified, and memory usage stayed within a
stable range instead of growing out of control. This does not yet prove that the
system is ready for Ethereum block proving in practice, but it suggests that
longer executions are not immediately blocked by catastrophic witness or memory
growth.

Therefore, the final answer is cautiously positive. The current implementation
does not yet provide a complete experimental proof that noise growth will remain
manageable at production scale, but the observed long-running behavior supports
the claim that the lattice-based design is viable and deserves further
optimization and larger-scale testing.
